% unidades/01-repaso-n4.tex
\chapter{\emoji{seedling} \ruby{復習}{ふくしゅう} — Repaso N4}
\unidadheader{01}{復習}{Repaso N4}{Consolidar las bases antes del nivel intermedio}

\begin{tcolorbox}[vocabbox, title={\emoji{pushpin} Vocabulario Clave}]
\begin{tabularx}{\linewidth}{llX}
\toprule
\textbf{Japonés} & \textbf{Español} & \textbf{Tipo} \\
\midrule
\vocabitem{復習する}{ふくしゅうする}{repasar / revisar}{v. suru}
\vocabitem{確認する}{かくにんする}{confirmar / verificar}{v. suru}
\vocabitem{練習する}{れんしゅうする}{practicar}{v. suru}
\vocabitem{間違える}{まちがえる}{equivocarse}{v. G2}
\vocabitem{覚える}{おぼえる}{memorizar}{v. G2}
\vocabitem{忘れる}{わすれる}{olvidar}{v. G2}
\vocabitem{理解する}{りかいする}{comprender}{v. suru}
\vocabitem{上達する}{じょうたつする}{mejorar (habilidad)}{v. suru}
\bottomrule
\end{tabularx}
\end{tcolorbox}

\subsection*{\emoji{speech-balloon} Frases Clave}

\frase{もう一度言ってください。}{Por favor dígalo otra vez.}
\frase{\ruby{間違}{まちが}えてもいいです。}{Está bien equivocarse.}
\frase{〜たことがあります。}{He hecho ~. (experiencia pasada)}
\frase{〜なければなりません。}{Tengo que hacer ~. (obligación)}
\frase{〜ようにしています。}{Intento / procuro hacer ~. (hábito intentado)}

\subsection*{\emoji{gear} Gramática — Repaso Rápido N4}

\patron{Estado continuo vs acción en progreso}
  {\ruby{動詞}{どうし}て\ruby{形}{けい} + いる}
  {\ej{\ruby{彼}{かれ}は\ruby{日本語}{にほんご}を\ruby{勉強}{べんきょう}している。}{Él está estudiando japonés. (en progreso)}
   \ej{\ruby{窓}{まど}が\ruby{開}{あ}いている。}{La ventana está abierta. (estado resultante)}}

\patron{Experiencia pasada}
  {〜た + ことがある}
  {\ej{すしを\ruby{食}{た}べたことがあります。}{He comido sushi alguna vez.}
   \ej{\ruby{富士山}{ふじさん}に\ruby{登}{のぼ}ったことがありません。}{Nunca he subido al Monte Fuji.}}

\patron{Obligación}
  {〜なければなりません / 〜なくてはいけません}
  {\ej{\ruby{毎日}{まいにち}\ruby{練習}{れんしゅう}しなければなりません。}{Tengo que practicar todos los días.}
   \ej{\ruby{薬}{くすり}を\ruby{飲}{の}まなくてはいけません。}{Tengo que tomar la medicina.}}

\comparacion{〜ている: dos usos}
  {acción en progreso\\\small\ruby{今}{いま}\ruby{食}{た}べている。\\\textit{Estoy comiendo ahora.}}
  {estado resultante\\\small\ruby{結婚}{けっこん}している。\\\textit{Está casado/a.}}
  {El contexto y tipo de verbo determinan el significado. Verbos de acción → progreso. Verbos de cambio de estado → estado resultante.}

\coloquial{〜なければなりません}{〜なきゃ / 〜なくちゃ}
  {Contracción muy común en habla informal. 行かなければなりません → 行かなきゃ。}

\culturabox{N4から先の日本語}{
  Al pasar del N5-N4 al N3, el japonés empieza a sentirse más \textit{natural}. Los patrones se vuelven más complejos, pero también más expresivos. La clave es leer y escuchar japonés real: noticias de NHK Easy, mangas, dramas. \ruby{継続}{けいぞく}は\ruby{力}{ちから}なり — La constancia es poder.
}

\lectura{\ruby{先生}{せんせい}からのメッセージ\\[0.4em]
  N4まで\ruby{勉強}{べんきょう}したみなさん、\ruby{本当}{ほんとう}によく\ruby{頑張}{がんば}りました。ここからはN3です。\ruby{文法}{ぶんぽう}がもっと\ruby{複雑}{ふくざつ}になりますが、\ruby{表現}{ひょうげん}も\ruby{豊}{ゆた}かになります。\ruby{毎日}{まいにち}少しずつ\ruby{練習}{れんしゅう}すれば、きっと\ruby{上達}{じょうたつ}できます。}
{Mensaje de la profesora\\[0.4em]
  A todos los que han estudiado hasta N4: realmente lo han hecho muy bien. A partir de aquí es N3. La gramática se vuelve más compleja, pero las expresiones también se enriquecen. Si practican un poco cada día, seguramente mejorarán.}

\begin{tcolorbox}[ejerciciobox, title={\emoji{pencil} Ejercicios de Repaso}]
\begin{enumerate}
  \item Conjuga en て-form: \ruby{書}{か}く、\ruby{食}{た}べる、する、くる
  \item Transforma: \ruby{毎日}{まいにち}\ruby{運動}{うんどう}する (→ obligación formal)
  \item ¿Cuál es la diferencia? \ruby{ドアが開いている}{ドアがあいている} vs \ruby{ドアを開けている}{ドアをあけている}
  \item Crea una oración con たことがある sobre una experiencia real tuya.
\end{enumerate}
\vspace{0.4em}
\textbf{Respuestas:}
1) \ruby{書}{か}いて・\ruby{食}{た}べて・して・きて \quad
2) \ruby{毎日}{まいにち}\ruby{運動}{うんどう}しなければなりません \quad
3) 開いている = la puerta está abierta (estado); 開けている = está abriendo la puerta (acción) \quad
4) (respuesta libre)
\end{tcolorbox}

\begin{tcolorbox}[kanjibox, title={\emoji{mahjong} Kanjis de la Unidad}]
\begin{tabularx}{\linewidth}{cllXX}
\toprule
\textbf{Kanji} & \textbf{On} & \textbf{Kun} & \textbf{Significado} & \textbf{Vocab} \\
\midrule
\kanjirow{復}{ふく}{—}{\ruby{repetir}{}/recuperar}{\ruby{復習}{ふくしゅう} / \ruby{復活}{ふっかつ}}
\kanjirow{習}{しゅう}{\ruby{なら}{う}}{\ruby{aprender}{}/practicar}{\ruby{復習}{ふくしゅう} / \ruby{習慣}{しゅうかん}}
\kanjirow{練}{れん}{\ruby{ね}{る}}{\ruby{entrenar}{}/refinar}{\ruby{練習}{れんしゅう} / \ruby{熟練}{じゅくれん}}
\kanjirow{確}{かく}{\ruby{たし}{か}}{\ruby{seguro}{}/confirmar}{\ruby{確認}{かくにん} / \ruby{確実}{かくじつ}}
\bottomrule
\end{tabularx}
\end{tcolorbox}
